\begin{table}[H]\centering
\caption{Факторное исследование Luna: число назначений и наблюдений.}
\label{tab:scale-quality}
\begin{tabular}{lrrrrr}\toprule
Условие & Назначено & Сгенерировано & Наблюдается & Успехи & Успех (\%)\\\midrule
D & 3000 & 2996 & 2951 & 1611 & 54,59\\
SN & 3000 & 2998 & 2953 & 1604 & 54,32\\
SR & 3000 & 2997 & 2952 & 1581 & 53,56\\
MN & 3000 & 2983 & 2938 & 1471 & 50,07\\
MR & 3000 & 2987 & 2942 & 1514 & 51,46\\
\bottomrule\end{tabular}
\par\smallskip\noindent «Наблюдается» означает доступные исходы качества на пригодных задачах при штатной проверке, а не отдельные задачи. Доля успеха равна числу успешных исходов, деленному на число наблюдаемых повторов. В каждом условии на задачу назначено три повтора.
\end{table}
